% ============================================================================
%  section/Conclusion.tex -- Conclusion
% ============================================================================
\chapter{CONCLUSION}

This minor project investigated and demonstrated the feasibility of a scalable
multi-hop LoRa network for intelligent forest monitoring, integrating
environmental sensing, edge-based machine learning, and low-power
communication into a single sensor node design.

The hardware prototype successfully collected temperature, humidity, and gas
concentration data and combined them through multisensory fusion to compute a
dynamic fire risk score. The testing results confirmed that combining multiple
environmental parameters provides a more reliable indication of fire risk than
individual sensor threshold methods and reduces false alarms caused by
temporary fluctuations in a single reading.

The TinyML inference pipeline, evaluated on the training dataset, demonstrated
that acoustic features could be effectively classified on-device, enabling the
detection of threats such as chainsaws, axe cuts, and gunshots close to where
the sound occurs. This edge-based approach avoids sending raw audio over the
network, significantly reducing bandwidth and energy consumption.

The LoRa communication simulation validated the ability of the network to
extend coverage beyond the range of a single gateway. The proposed LDSE
multi-hop protocol is expected to provide an energy-efficient, slot-scheduled
routing scheme that maintains synchronization across the network while
allowing nodes to remain in deep sleep outside their assigned time slots.

Overall, the design demonstrates that a low-cost, scalable, and
energy-efficient solution based on LoRa and TinyML is achievable for real-time
forest monitoring in community forests and protected areas. The remaining
tasks described in the preceding chapter define the path toward a complete
field-deployable implementation.